One property of the Lotka--Volterra benchmark bounds what its statistical numbers can mean, and we state it here rather than leave it to be discovered. The simulator has an absorbing state: if either population reaches zero before $T_{\max}$ the run returns a fallback discrepancy, which a smooth kernel weights to effectively zero. At the configured parameters that is the usual outcome. Survival is $2.6\%$ of $2000$ simulations at the true parameter, and $1.7\%$ over $400$ uniform prior draws, of which $92.5\%$ never survived one of $25$ attempts. The observed dataset is itself survival-conditioned --- it is re-drawn until a trajectory survives, which here takes more than eleven attempts --- so the target is $p(\theta \mid s_{\mathrm{obs}}, \text{survived})$ against a survival-conditioned $s_{\mathrm{obs}}$, which is internally consistent: the extinct particles' zero weight \emph{is} the $P(\text{survive}\mid\theta)$ factor of the likelihood and belongs there. The consequences split. Throughput and scaling on this benchmark are unaffected, because an extinct simulation costs what a surviving one costs and that cost is the systems quantity under test. Posterior recovery is a different matter: the informative sample is ${\approx}2\%$ of the nominal budget, so a $1.2\times10^6$-evaluation run carries the information of ${\approx}25{,}000$, and any method comparison on it is substantially a comparison of how each copes with a $98\%$-rejection regime. We therefore report Lotka--Volterra as a systems benchmark and treat its posterior curves as we treat Cellular Potts's: as diagnostics that bound no method's attainable accuracy. Computational metrics: simulation throughput, worker utilization / idle fraction, and wall-clock time to a fixed posterior-error target. Convergence curves are placed on a shared time grid via last-observation-carried-forward resampling so asynchronous (fine-grained) and synchronous (per-generation) records compare fairly.
